\section{Comparison with the ordinary factor categories}
\label{sec:comparison}

Proposition~\ref{prop:curvature} places the immersion and the embedded
test branes in a proper uncurved cyclic category.  To use the ordinary
mirror equivalences, we must identify the category generated by the
embedded tests inside that category.  The needed comparison concerns
only words of distinct, ordered copies.  This restriction gives isolated
asymptotic intersections at every end, including every internal end of a
broken polygon, and removes both self-intersection sectors and unit
conventions from the analytic comparison.

Fix one of the finite full cyclic categories constructed in
Proposition~\ref{prop:curvature}, and select its embedded copy labels.
Choose any total ordering of these labels and write the copies in that
order as
\[
 K_a=\Sigma_a\times G_a,\qquad 1\leq a\leq N,
\]
where each is a copy of a test from Proposition~\ref{prop:realization}.
Omit $\Sigma_a$ when $Y$ is a point.  Here $\Sigma_a$ and $G_a$ include the Hamiltonian
perturbations and transported brane structures of the copies.  As in
the preceding construction, all pairs are transverse and there are no
triple intersections, also in each factor separately.  The fixed split
almost complex structure $J_0$ has no nonconstant spheres, and none of
these embedded copies bounds a nonconstant $J_0$-holomorphic disk.
A \emph{strict word} means a sequence of labels
$a_0<\cdots<a_k$.  Its inputs lie in the intersection complexes for
$(K_{a_{i-1}},K_{a_i})$, and its output lies in the complex for
$(K_{a_0},K_{a_k})$.  We initially omit identities and self-morphism
spaces.  Thus the directed category in this section means precisely
these complexes and their operations on strict words.

\begin{proposition}[Directed comparison]
\label{prop:directed-comparison}\label{prop:comparison}
For each fixed finite full cyclic category above and every total
ordering of its embedded copy labels, its strict directed sector is
quasi-isomorphic, after
an extension of its Novikov field if necessary, to the corresponding
ordinary directed Fukaya category of the embedded product copies.
When $Y$ is present, this ordinary directed category is
quasi-isomorphic to the directed restriction of the $A_\infty$ tensor
product of the two ordinary factor categories.  These comparisons use
the specified brane labels and their usual graded intersection
complexes.  Consequently the ordinary mirror functors in the two
factors give a comparison with the directed category on the specified
objects $Q_\ell\boxtimes N_M$.
All the finitely many ordering comparisons for this same full cyclic
category may be made over one field extension.
\end{proposition}

The two comparisons change different parts of the construction.
For the first, we keep $J_0$ and the intersection generators fixed
and change disk integration into regular Floer counts by introducing
Hamiltonian terms.  Section~\ref{comp:virtual} verifies the relative
chart and gluing hypotheses for that change.  It then removes the
energy cutoffs algebraically, while retaining the directed sector
of the originally fixed full cyclic category.

For the second comparison, we retain the regular counts and change
their domain parameters.  The product category uses the same polygon
in both factors; the tensor product uses an associahedral diagonal.
Sections~\ref{comp:domain-chains}--\ref{comp:orientations} construct
and count a homotopy between these choices, including the boundary
and sign identities that make the counts an operad map.
Section~\ref{comp:operadic-completion} then identifies the result with
the tensor product of the ordinary factor categories.  All these
comparisons concern embedded copies, and ordinary mirror symmetry
is used in the individual factors.  The freedom to order the same
copies in every way will allow us to test compositions through three
distinct copies and prove that the copy arrows are invertible in
Lemma~\ref{realization:copies}.

\subsection{Regular Hamiltonian models with fixed intersections}
\label{comp:regular-models}

For each ordered pair choose strip data of the form $H(t,x)\,dt$,
$0\leq t\leq1$, with $H$ small and identically zero on fixed small
balls around the intersections of that pair.  Take $H$ sufficiently
small on the complement as well.  Its chords are then exactly the
original intersections.  On a stable polygon, with at least two
inputs, use a domain-dependent Hamiltonian one-form whose tangential
restriction to every Lagrangian boundary arc is zero.  The almost
complex structure remains $J_0$.  Product data are the sums of the
corresponding data in the two factors.

We require agreement with the pair data on the terminal parts of the
external ends and on long necks.  On polygon cores the data are smooth
up to domain faces and restrict there to the data on the component
polygons.  More precisely, use the usual consistent strip-like ends
and simultaneous plumbing collars on the compactified domain spaces.
If $T$ is a neck length, smoothness of the core coefficients is taken
in a coordinate $\exp(-cT)$, for a fixed $c>0$.  Thus their difference
from the limiting component coefficients, together with every
required length derivative, is exponentially small.

These conditions leave enough freedom for regularity.  To see their
compatibility, first choose data agreeing with their component data
on smaller collars and carrying the prescribed strip model near each
node.  Retain that strip model on terminal end and neck
neighborhoods; variations occur on the complementary core patches.
The restrictions to intersecting domain faces agree, so extend their
coefficients by the ordinary extension theorem for functions on a
manifold with corners, using inclusion--exclusion and a partition of
unity.  The zero tangential boundary condition is linear and is
preserved by this extension.  At later arities one may shrink the
fixed neighborhoods, since the lower-arity data already equal the
strip model there.  This construction gives smooth extensions with
equality on faces; it does not impose equality with component data on
an entire domain collar.  In particular, at any smooth domain, even
one lying in a plumbing collar, there remain open interior core
patches on which the Hamiltonian data may vary.  The radial
plumbing description and strip height used here are also those of
\cite[Proposition~7.5]{A8}.

\begin{lemma}
\label{comp:regularity}
For the finite list above, the Hamiltonian data can be chosen small,
with the stated face and strip conditions, so that the ordinary
strict-word operations are regular.  They can also be chosen regular
over any prescribed countable collection of smooth parameter chains
in products of domain spaces, including all their relative face
interiors and all component configurations at those faces.
Compactness below an area bound has only polygon and strip
degenerations.  All nonconstant strips and all contributing stable
polygons have positive area, with a common positive lower bound for
this finite directed system.
\end{lemma}

\begin{proof}
The energy identity compares analytic energy with ordinary symplectic
area and a curvature term supported on the cores.  The integral of
that curvature is uniformly bounded for each word, because the pure
strip model has zero curvature on every neck.  There is no boundary
term along a labeled arc, and the chords have not moved.  On strips,
analytic energy equals area.  Areas are additive at broken
configurations and are constant in a relative class during the
allowed homotopies of data.

At a small-scale bubble the Hamiltonian terms disappear, leaving a
$J_0$-holomorphic sphere or a disk on one of the embedded copies.
Neither exists.  Compactness therefore reduces to ordinary Floer
compactness with isolated transverse asymptotes.  Below a fixed area
bound, only finitely many relative classes are realized, even before
regularity: near a limiting broken configuration the total relative
class is fixed.  For regular data this gives finite zero-dimensional
counts and compactifications of the one-dimensional moduli spaces.

Nonconstant strips have positive area and a positive energy gap.
Stable polygons also have a uniform positive area gap once the data
are sufficiently small.  Otherwise a sequence of data tending to
zero and polygons with areas tending to a nonpositive value would
have analytic energies tending to zero.  Compactness, including on
long necks, would make every component constant.  A stable vertex in
a strict word has at least three distinct boundary labels, so such
a constant component would give a triple intersection.  This is
excluded.  There are only finitely many strict words, which gives a
single small-data neighborhood and a common gap.

For strips, regularity follows by varying $H(t,x)$.  The regular-point
argument based on unique continuation gives dense injective points
of $(t,u(s,t))$ on every nonconstant trajectory, away from its
asymptotes; it applies with Lagrangian boundary just as in the Floer
strip argument of \cite{FloerHoferSalamon}, with the corrected
linear-algebra step in \cite{FloerHoferSalamonErratum}.
The required variations
may be supported away from the intersection balls.  Indeed a
nonconstant strip cannot remain entirely in those balls: the local
sheets are exact there, so its area would be zero.  A constant strip
has an invertible transverse operator.

For polygons, domain-dependent Hamiltonian variations on interior
core patches give transversality without a simplicity assumption on
the map.  The same argument applies over any fixed smooth parameter
chain.  On a domain face, the vertices in one factor have different
tuples of boundary labels, by strictness.  Their variations can
therefore be chosen independently.  A variation on the relative
interior of a face extends to the full system by the face-extension
construction above; the restrictions on higher families do not
constrain variations on these component core patches.  Thus the
universal linearized problem is surjective on every relevant
stratum.  Sard--Smale, with countable exhaustions by energy, class,
and parameter stratum, supplies simultaneous regularity.  Applied to
the diagonal chains it gives regular product polygon operations.
Product strip regularity is the direct-sum statement for the factor
strip operators.
\end{proof}

\subsection{Relation to the ordinary Fukaya categories}
\label{comp:ordinary}

The preceding data compute the ordinary factor categories used in
mirror symmetry.  For the graphs this is ordinary weakly exact
Floer theory.  For the quartic spheres, take their disk-excluding
almost complex structure to be $J_0$, as permitted in Seidel's
construction \cite[Section~8]{SeidelQuartic}.  It is also sphere-free,
so regularity at these fixed data needs no additional bubble-incidence
condition.  Extend the strict-word data to all ordinary words on the
same copies; repeated labels use nondegenerate chord perturbations.
The extension is again inductive over domain faces, with the same
terminal strip models and smooth core coefficients.

We explain why the slightly more flexible collar condition used
above gives the same ordinary category.  One can put two choices on
separately labeled copies and extend their mixed-word data.  Energy
and compactness for the comparison follow from the same no-bubbling
argument.  At a broken curve the core coefficients differ from their
limits exponentially in the neck lengths, so ordinary gluing with
isolated asymptotes applies.  Continuation, binary products, and
self-Floer units then identify the two labels of each object.  This
is the ordinary comparison argument; the collar proof in
\cite[Proposition~7.5]{A8} uses exactly these hypotheses and even
allows decaying errors in strip-end coefficients.

When comparing with an ordinary choice that has varying almost
complex data, the quartic factor requires the usual low-dimensional
sphere-avoidance argument.  The prescribed disk-excluding structure
at each boundary label is retained, which excludes boundary bubbles.
In real dimension four with $c_1=0$, a simple unmarked sphere has
expected dimension $-2$.  Adding the two incidence marks, one on
the sphere and one on the principal component, and imposing the
four-dimensional matching condition leaves a loss of two dimensions.
Removing a sphere tree does not change the principal component
index.  Hence sphere incidence can be avoided in the zero- and
one-dimensional systems used for operations and comparisons.  For
a multiply covered sphere, its image lies on a simple sphere, so the
same incidence exclusion suffices.

The required transversality uses almost complex variations for the
sphere and Hamiltonian variations for the principal component,
including evaluation at the incidence point.  Constant or contracted
components do not alter this test.  In a relative construction,
previously fixed faces already satisfy the exclusion, and the new
variations are available on the relative interiors.  The terminal
strip and neck data may be chosen sphere-free; strip data depending
on strip height form only a one-parameter family.  Continuation data
vary on compact cores, where the same argument applies.  These are
the sphere-exclusion and comparison mechanisms of
\cite[Section~8]{SeidelQuartic}.  Thus the mixed-copy construction
identifies our regular systems with conventional ordinary systems,
including cohomological units.  The actual Hamiltonian sphere copies
represent the original objects by
\cite[Lemma~8.14]{SeidelQuartic}.

If one uses capped-generator Novikov exponents in the rational-brane
normalization, rather than ordinary polygon areas, the coefficients
differ by fixed input and output scalar factors.  Rescaling the
generator lines after the permitted scalar extension makes these
descriptions agree.  This completes the identification of the
regular factor models.  It remains to compare them with the disk
integration operations that were used to include the immersion.

\subsection{The relative comparison from disk integration}
\label{comp:virtual}

Join the zero Hamiltonian data to the regular data by a small path,
constant near its endpoints.  Keep $J_0$ and every intersection chord
fixed.  We use the disk Kuranishi and CF-perturbation construction
over this path.  Here a CF perturbation is the compatible finite
dimensional perturbation used to integrate along a Kuranishi
correspondence.  The comparison will use its relative extension and
Stokes properties; it does not require cyclic perturbations along
the path.

\begin{lemma}
\label{comp:pseudoisotopy}
For every finite energy and arity cutoff, the strict-word disk
integration operations at the zero end and the ordinary regular
counts at the other end are related by a filtered pseudo-isotopy.
The linear comparison has the identity reduction on the
intersection generators.  The data at either endpoint may be
recovered up to the usual finite-cutoff pseudo-isotopy.
\end{lemma}

\begin{proof}
We verify that the relative disk construction applies to the
Hamiltonian terms just introduced.  The analytic point is that the
unchanged transverse asymptotes and the controlled core coefficients
give uniform gluing estimates at every corner.  With these estimates,
the cited chart and relative perturbation constructions apply with
the path parameter included.  The inputs are the disk charts
from obstruction bundle data, their compatible products on normalized
corners, and the relative CF extension and Stokes theorems used in
the cyclic construction
\cite[Sections~9--10 and~12]{AFOOO2026}.  For the chart and gluing
theorems, we use \cite[Theorem~7.1, Condition~7.11, and
Lemma~7.12]{FOOOChartsI}, the disk-component-wise compatibility
construction of \cite{FOOOChartsII}, and the exponential gluing
estimates in \cite[Assumption~3.12 and Theorems~3.13, 6.4,
8.16--8.17]{FOOOGluing}.  The following verifications explain their
application to this relative problem.

\paragraph{Charts and the path parameter.}
Use stabilized strip and polygon domains, with the path coordinate
included among the parameters.  The Hamiltonian terms on a core are
smooth lower-order terms in the Cauchy--Riemann equation.  Compact-core
obstruction sections, transported and summed as in the disk-chart
construction, make the linearized equation surjective at a fixed
domain and fixed path parameter.  Interior test sections suffice:
the adjoint unique-continuation argument is unchanged by these
smooth lower-order terms.  This also makes the path projection
submersive.  Strip translation is removed by the usual local slice.
When stabilization uses auxiliary domain points constrained to
target slices, suitable immersive points exist: the maps under
consideration have strictly positive area.  Strict words give no
map isotropy, and the no-bubbling argument requires no sphere chart.

\paragraph{Coefficients near broken configurations.}
At a gluing seam the ends lie in the balls where the Hamiltonian
vanishes.  The seam equation therefore coincides with the
unperturbed transverse-end equation.  A nonconstant strip core down
a neck may carry a Hamiltonian term, but that term is translation
invariant and transports smoothly with the core.

There is a coordinate issue when a domain neck contains several
Floer breaking seams.  Its domain length is the total length of the
intervening pieces.  Consequently its exponential plumbing
coordinate is a product of the exponential seam coordinates, times
a smooth positive factor.  The core coefficients retain their
smoothness and exponential derivative estimates after this
contraction.  Equivalently, in stabilized charts keep the auxiliary
marks that locate all the retained cores.

For completeness, use radial variables
\[
 r_l=e^{-\pi s},\qquad r_r=e^{-\pi(T-s)}
\]
on a boundary neck, together with the strip height.  Transport to a
stabilized strip core multiplies each one-sided radial variable by
the intervening radial factors and smooth core factors.  Changes of
boundary-node coordinate have a positive nonzero leading real
coefficient.  An origin shift of order $r_lr_r$, after division by
one radial variable, still contains the opposite radial factor.
The strip coframe thus tends smoothly to the limiting strip
coframe in these coordinates.  These are the radial collar
calculations in \cite[proof of Proposition~7.5]{A8}.
For coordinate changes between map charts one retains the marks
needed for local stability, including marks on strip components.
There is no requirement to express the inverse total neck length
smoothly in the inverse lengths of a contracted chain.

\paragraph{Uniform gluing estimates.}
Choose exponential weights below the spectral gaps of the finitely
many transverse asymptotic operators.  On a joined strip of length
$T$, a smoothed version of
$e^{\delta\min(s,T-s)}$ gives a suitable weight.  Augment the
component operators by kernel observables on compact cores and by
the obstruction variables.  Pregluing errors are exponentially
small in a slightly weaker weight: the maps decay at the ends, the
seam equations agree, and the core coefficient errors are
exponentially small.

To build a parametrix, split the forcing near each seam midpoint in
a fixed-width overlap, solve with the augmented component inverses,
and splice over intervals whose lengths are proportional to $T$.
The component one-sided weights dominate the joined weight on their
respective sides, uniformly on that overlap.  The parametrix error
is the sum of an $O(T^{-1})$ cutoff error and an exponentially small
coefficient error.  There is no matching zero mode at a transverse
intersection.  A Neumann-series correction therefore gives a
uniform augmented inverse; the Fredholm indices agree under
gluing.  The quadratic estimate is uniform as well, since all new
lower-order terms occur on fixed compact cores and the seam
equation near the joined values is unchanged.

The differentiated estimates have the same form.  Core length
derivatives of coefficients are exponentially small.  Derivatives
of middle cutoffs and reparametrizations multiply exponentially
decaying map terms.  Passing to weaker weights accommodates every
fixed finite number of derivatives, for several necks and with
some nodes still broken.  Hence the usual disk-chart smoothness in
inverse-length corner coordinates applies.  These are precisely
the exponential estimates required by the cited FOOO gluing
theorems.  Alternatively, the isolated-end case of
\cite[Lemma~4.11 and Proposition~4.12]{A8} uses the same augmented
surjectivity, spectral gaps, core transports, and exponentially
controlled coefficients.  Its strip splice estimates are estimates
for finite-rank systems and apply in the present dimension.  The
converse gluing assertion follows from closeness on the cores and
decay along the joining ends, giving the local chart coverage.

\paragraph{Compatible charts and relative perturbations.}
At a normalized corner, choose the product of the component
charts over their common path parameter, summing their
core-supported obstruction data.  At every iterated corner the
limiting equation data agree by the face construction.  The
disk-component-wise compatibility procedure now uses the same
stabilizations, nested neighborhoods, transported core sections,
openness of augmented surjectivity, and obstruction-data
enlargements as for unforced disks.  Compactness gives a finite
construction at each cutoff.  The preceding estimates supply the
inverse-function and gluing steps, with the path parameter included.
Thus the smooth lower-order core terms introduce no new condition
in that construction.  There are no forgetful-map or cyclic
compatibility requirements on strict words, and all seam matches
are at points.

Use the usual outer collars and relative CF extension, preserving
the component products and the submersive path projection.  Stokes
then gives the pseudo-isotopy equations with the required face
products.  Endpoint chart choices need only be recovered up to a
finite-cutoff pseudo-isotopy, as in the ordinary relative
construction.  At the regular endpoint, a virtual zero-dimensional
moduli space has no boundary strata and is already an oriented
zero-manifold without isotropy.  Its Kuranishi integral is its
ordinary signed count, also if one uses a chart with redundant
obstruction directions.  Thus the endpoint operations are the
ordinary operations from Lemma~\ref{comp:regularity}.

The positive-energy gap makes all nonidentity corrections to the
linear comparison have positive valuation.  Its reduction is
therefore the identity on the unchanged intersection generators.
\end{proof}

All orientations in this comparison come from the same graded spin
determinant lines, with zero background class.  The comparison of
sign conventions is proved in Section~\ref{comp:orientations} below.
If an ordinary convention reverses arrows, use the opposite
category and the corresponding algebraic duality; no reflection of
holomorphic disks is involved.

We next remove the cutoffs without imposing simultaneous geometric
choices at all orders.  Let $\Lambda$ be the Novikov field in use,
with valuation $v$.  There are finitely many strict words in the
finite list, and every intersection complex is finite dimensional.
Thus the equations for an $A_\infty$ isomorphism between its two
directed structures are a finite polynomial system over $\Lambda$.
Include variables for the inverse of its linear part.  For any
valuation bound, Lemma~\ref{comp:pseudoisotopy}, together with the
finite-cutoff comparisons of the chosen promoted cyclic model, provides an
approximate solution.  Its coordinates have nonnegative valuation,
as do the inverse linear-part coordinates, because that linear
part has invertible reduction.  In particular, the source of
these comparisons is the restriction of the fixed full cyclic category.
Changing chart or perturbation choices for a particular ordering is
accounted for by these endpoint pseudo-isotopies; it does not replace
that ambient category.

If this finite system had no solution over an algebraic closure of
$\Lambda$, the Nullstellensatz would give a fixed polynomial
combination of its equations equal to $1$.  Evaluation at an
approximate solution of sufficiently large valuation contradicts
this identity: all polynomial coefficients have a fixed lower
valuation bound, whereas the equation errors have arbitrarily high
valuation.  A solution therefore exists after field extension.
The argument applies to every ordering of the same finite label set.
There are only finitely many orderings, so all their polynomial
solutions may be chosen over one algebraically closed extension of
the original coefficient field.  This proves the first assertion of
Proposition~\ref{prop:directed-comparison}.  We now retain the
regular ordinary models and compare their product operations.

\subsection{Chains of polygon domains}
\label{comp:domain-chains}

Write $\mathcal K_j$ for the compactified space of disk domains with $j$
inputs and one output, $j\geq2$.  It has dimension $j-2$ and its
faces are products of lower polygon spaces.  We use its usual
smooth face structure from strip-like ends and plumbing.  A pair
of factor polygons has domain parameter in $\mathcal K_j\times \mathcal K_j$.
The product Fukaya operation uses the diagonal copy of $\mathcal K_j$ in
this product.  Replacing this diagonal by a cellular operad
diagonal will give the tensor product.

We need a chain model in which this replacement can be counted
geometrically.  A geometric chain in $\mathcal K_j^2$ is a smooth map
\[
 c:P\longrightarrow \mathcal K_j^2
\]
from a compact oriented smooth manifold with corners and finitely
many faces.  Every relative face interior of $P$, including its
interior, must map into a relative face interior of $\mathcal K_j^2$.
Impose also the following condition on the way target faces are
left.  If a face of $P$ has codimension one in a containing face,
the inward derivative of $c$ is strictly inward at every target
face hypersurface that contains the smaller image interior but
does not contain the larger image interior.  Derivatives here are
taken in the exponential plumbing coordinates.  The condition
allows several target nodes to be smoothed by one inward
parameter.

Take linear combinations of these maps, with the usual orientation,
reparametrization, and disjoint-union relations; the boundary is
the sum over oriented codimension-one faces.  A map that ignores
a direct product factor of positive dimension is declared zero.
This is a subcomplex relation: if the ignored factor has dimension
one, its signed endpoints cancel, and in larger dimensions its
boundary factors are still ignored.  Products, face inclusions,
diagonals, and grafting in both target coordinates preserve the
face condition.  With the identity in arity one adjoined, these
chains form a nonsymmetric differential graded operad, denoted
$\mathcal G$; a chain of dimension $d$ has cohomological degree
$-d$.

\begin{lemma}
\label{comp:chain-acyclic}
In every arity $j\geq2$, the chain complex $\mathcal G(j)$ has the
homology of a point, with the ordinary augmentation in degree zero.
\end{lemma}

\begin{proof}
Choose a smooth homotopy of $\mathcal K_j^2$ that immediately moves every
boundary point into the interior, with strictly inward initial
derivative at all its incident hypersurfaces, and then contracts
within the interior.  Such a homotopy is obtained from an inward
vector field followed by a contraction of a compact subset of the
open cell.  Apply it to a chain and use its prism.

At time zero the strict inward condition gives exactly the
required condition on the new prism faces.  At positive time the
image is in the target interior, so no further inward condition
remains.  The prism identity is therefore valid in $\mathcal G$.
At the constant endpoint every positive-dimensional chain ignores
its whole domain as a product factor and is zero.  A
zero-dimensional chain is reduced to its augmentation times the
chosen interior point.  This gives the asserted contraction.
\end{proof}

We can now choose the domain homotopy, before choosing the Floer
data with which it will be counted.  Let
$\mathcal A_\infty=C_{\mathrm{cell}}(\mathcal K_\bullet)$ be the cellular
dg operad of associahedra, with the same cohomological convention.
Choose a cellular diagonal
\begin{equation}
 \Delta:\mathcal A_\infty\longrightarrow
             \mathcal A_\infty\otimes\mathcal A_\infty
 \label{comp:cellular-diagonal}
\end{equation}
lifting multiplication.  Such a diagonal exists by induction on
arity, or by the free resolution description; see
\cite{MarklShnider}.  No coassociative choice is required.

There are two maps from $\mathcal A_\infty$ to $\mathcal G$.
The first takes the fundamental chain of $\mathcal K_j$ to its actual
diagonal in $\mathcal K_j^2$.  The second applies $\Delta$ and then the
cross product of geometric fundamental chains.  Both preserve
the augmentation.  Lemma~\ref{comp:chain-acyclic} and cofibrancy
of the $A_\infty$ operad imply that they are dg-operad homotopic.
One can construct the homotopy arity by arity, using polynomial
differential forms on an interval as a path object: each
obstruction is an augmentation-zero cycle in an acyclic arity.

Fix these two maps and one such homotopy in $\mathcal G$.
Their coefficients involve only countably many geometric chains.
Take the dg suboperad $\mathcal G_0\subset\mathcal G$ generated
by those chains, and include the representatives and chain
relations used in this construction among the choices below.
The countable family remains countable after adjoining products,
faces, and maps with ignored factors removed.  If a relation requires
additional representatives, add finite witnesses for a countable
basis of the relation space and repeat this closure.  This still
uses only countably many maps.  Choose regular
split Hamiltonian data for this family by
Lemma~\ref{comp:regularity}.  We require counts only on
$\mathcal G_0$; no assertion of its acyclicity is needed.

When a parameter lies in a
target face, count the polygons at all its vertices, matching the
internal ends and summing over their discrete intersection
generators.  Weight a configuration by its total area.  The
fixed-domain Fredholm index is
\[
 |x_0|-\sum_{i=1}^j|x_i|.
\]
Consequently a rigid count over a $d$-dimensional parameter chain
acts in degree $-d$.

\subsection{Compactification, gluing, and the chain identities}
\label{comp:chain-counts}

We show that these parametric counts, with the orientation
normalization specified in the next subsection, define a dg
operad map from $\mathcal G_0$ to the endomorphism operations of the directed
intersection complexes.  This is the geometric assertion needed
for the diagonal replacement.  It consists of compatibility with
the ignored-factor relation, boundary, and grafting.

If a chain ignores a positive-dimensional product factor, a
zero-dimensional count on it would have negative expected
dimension after that factor is removed.  Include the maps with
those factors removed among the regularity requirements.  Their
moduli spaces are empty, so the ignored-factor relation is
respected.

Consider next a one-dimensional parametric moduli space below a
fixed area bound.  Compactness in
Lemma~\ref{comp:regularity} permits two phenomena: approach to a
face of $P$ and breaking of a Floer strip.  Passing to a
codimension-$r$ face of $P$ lowers the virtual dimension by $r$.
This remains true when the face creates several nodes in one or
both polygon factors.  Indeed the fixed-domain indices add under
matching at isolated transverse intersections; the domain
parameters are supplied by $P$.  Each additional nonconstant
quotient strip lowers dimension by one more.  Stratum regularity
therefore excludes a strip breaking over a parameter boundary in
a one-dimensional problem, including a breaking down any newly
formed neck.  The same argument excludes boundary strata in a
zero-dimensional problem.  Repeated labels on strip components
cause no additional fiber-product transversality requirement,
since every match is discrete.

The uniform curvature bounds apply while a chain parameter tends
to a face.  Thus this is compactness for the parametric problem,
not merely compactness at each fixed parameter.  Its limiting
configurations at a target face are exactly the component
polygons and the strip breakings just described.

We must also know that each parameter-face configuration is the
end of the one-dimensional moduli space with its prescribed
orientation.  Use a local collar $(t,y)$ on $P$, with $t\geq0$
increasing inward.  For every newly smoothed neck its plumbing
coordinate has the form
\begin{equation}
 \lambda_i(t,y)=t\,c_i(t,y),\qquad c_i(t,y)>0,
 \label{comp:simultaneous-plumbing}
\end{equation}
by the face condition.  The neck lengths tend to infinity, and
their derivatives in $y$ stay bounded.  At $t=0$, regularity of
the nodal problem includes the $y$ variables.  Its augmented
linear operator is therefore the invertible operator needed for
the rigid parametric configuration.

For fixed small $t>0$, simultaneously preglue all these necks and
apply the weighted inverse construction from
Lemma~\ref{comp:pseudoisotopy}, including the $y$ variables in the
linearization.  Core coefficients converge to their limiting
values.  Differentiating in $y$ introduces only bounded
length derivatives multiplied by core-supported terms or by
terms decaying along the necks.  On the component halves, the
length-dependent transport has the same exponential estimates.
The implicit function theorem gives a unique glued solution for
each sufficiently small $t>0$, and ordinary gluing coverage
identifies all solutions near the nodal configuration with this
end.  Transport of the determinant orientation gives the boundary
orientation, with $t$ increasing inward.

This argument uses oriented-end coverage and the inverse estimates
for $t>0$.  It does not assert that the solution map is smooth at
$t=0$ in exponential plumbing coordinates.  The disk charts used
in Lemma~\ref{comp:pseudoisotopy} have the usual inverse-length
smoothness; these two uses of gluing require different conclusions.

At an edge already internal in a target face, a broken strip can
occur as an output breaking of the child polygon or as an input
breaking of the parent.  These are the same matched
configurations with opposite boundary signs, so they cancel.
External strip breakings give the Hom differential.  Grafting of
chains gives composition by the same discrete matching rule.
Finally, on a product of chains in the individual domain spaces,
the count factors: regularity makes a rigid product configuration
a product of rigid configurations in the two factors.  The
remaining point is that all these identities use a single
consistent sign convention.

\subsection{Determinant lines and signs}
\label{comp:orientations}

We use the graded orientation lines of Floer generators and the
spin determinant capping and gluing theorem, with zero background
class; for the immersed conventions see
\cite[Sections~4--5]{AkahoJoyce}.  Our factor Lagrangian dimensions
are even.  This parity removes the extra signs that would come
from moving the oriented rank contribution of a capped disk.

For a generator $x$, choose an input cap of index $|x|$ and an
output cap of reciprocal index.  Their parities agree in even
dimension.  Let $o_x$ be the input-cap line.  Order a fixed map
operator first, its input caps in order next, and its output cap
last.  The spin orientation of the resulting capped disk
identifies the map determinant symbol with the graded Hom line
\[
 \Hom\bigl(o_{x_1}\otimes\cdots\otimes o_{x_j},o_{x_0}\bigr).
\]
Pair reciprocal cap lines so that the invertible transverse
operator on a constant strip gives the identity of $o_x$.

The determinant gluing comparison at a specified end, with its
length held fixed, is local at that end after these caps are
inserted.  Determinant excision detaches the outer framing and
disk contributions; sphere determinants are complex oriented.
Only the two node caps remain in the comparison.  We record the
permutation that recovers the usual Hom sign.  Let a child map
determinant have parity $b$, let the sum of the preceding parent
input-cap parities be $l$, and let the node cap have parity $q$.
The sum of the child input-cap parities is $b+q$.  Starting with
the parent and child determinants, the external input caps, the
parent output cap, and the reciprocal node pair, reorder so that
each determinant is followed by its own caps.  The permutation
exponent is
\[
 bl+q.
\]
For example, if $r$ is the sum of trailing parent input-cap
parities and $o$ is its output-cap parity, the exponent before
simplification is
\[
 b(l+q+r+o)+q(b+q+r+o)+(r+o)(b+q)=bl+q
 \quad\text{in }\Z/2.
\]
The term depending only on the node, including $q$, is absorbed by
the reciprocal contraction already normalized on the constant
strip.  Fixed-operator gluing therefore gives Hom composition
with Koszul exponent $bl$.  This calculation applies to formal
Fredholm capping and gluing, independently of the existence of a
rigid holomorphic curve.

For direct sums choose the product cap to be the direct sum of the
factor caps, with line $o_x\otimes o_y$.  In the order
$o_x\otimes o_y\otimes o_x^\vee\otimes o_y^\vee$, its contraction
first moves $o_y$ past $o_x^\vee$ and then applies the two factor
contractions.  The crossing contributes $(-1)^{|x||y|}$, exactly
the graded tensor-evaluation sign.  The constant product strip is
the direct sum of the constant factor strips, so its determinant
unit is the tensor of their units and fixes this same reciprocal
normalization.  At a node the factors
$(-1)^{q_x}(-1)^{q_y}=(-1)^{q_x+q_y}$ absorb the product node parity;
the intervening mixed permutations remain precisely the tensor
Hom signs.  Determinant excision and product spin frames give
this direct-sum identification, with no additional rank sign in
our even dimensions.

For a parameter chain $c$ of dimension $d$, order its orientation
before the map determinant and write $f(c)$ for the resulting raw
count.  Composing chains of dimensions $p$ and $q$ introduces
$(-1)^{pq}$: the second parameter passes the first map
determinant.  The same sign appears for cross products.  Define
\begin{equation}
 \gamma(c)=(-1)^{d(d-1)/2}f(c).
 \label{comp:count-normalization}
\end{equation}
The identity
\[
 \frac{(-1)^{(p+q)(p+q-1)/2}(-1)^{pq}}
      {(-1)^{p(p-1)/2}(-1)^{q(q-1)/2}}=1
\]
shows that $\gamma$ respects both composition and cross products.

For the boundary identity let $v$ be the raw strip map with a fixed
translation convention.  An output strip places its map
determinant before the polygon determinant, so its translation
direction passes the $d$ parameter directions.  Write $s\in\{1,-1\}$
for its comparison with outward-normal-first orientation before
that permutation.  At an input the strip translation has the
opposite end direction, and its determinant follows the polygon
determinant.  Moving this translation past both the parameter and
the rigid polygon determinant adds no further parity.  Thus the
oriented boundary equation is
\begin{equation}
 f(\partial c)+s(-1)^d v\circ f(c)
       -s\sum_i f(c)\circ_i v=0,
 \label{comp:raw-boundary}
\end{equation}
where the slot compositions already include their graded Hom
signs.  On product generators $v$ has the tensor differential
rule.

The same calculation verifies internal-edge cancellation.  Put
the internal strip determinant in the same position on both
descriptions of the broken configuration.  Every permutation used
to do this is also the permutation in Hom composition, by the
fixed-operator calculation above.  The two translation directions
are opposite, which gives cancellation.  At a parameter face,
determinant gluing is compared at fixed positive smoothing values.
The parameter orientations then compare by their ordinary
boundary orientation, with $t$ inward in
Equation~\eqref{comp:simultaneous-plumbing}.  There is no separate
normal sign for each node smoothed by that same $t$.

Set the Floer differential to be $s v$.  Since $f(c)$ has degree
$-d$, Equation~\eqref{comp:raw-boundary} gives
\[
 d_{\mathrm{End}}f(c)=(-1)^{d+1}f(\partial c).
\]
The quotient of the normalization signs in dimensions $d$ and
$d-1$ is $(-1)^{d-1}$.  Therefore
\begin{equation}
 d_{\mathrm{End}}\gamma(c)=\gamma(\partial c).
 \label{comp:dg-count}
\end{equation}
Together with the compactification and grafting arguments, this
proves the promised dg operad count identities.

We finally compare this convention with conventional Fukaya
signs.  Match the cap lines and the order of geometric inputs,
using the same dg or suspended translation throughout.  The
remaining difference is removed by strict sign changes at
generators, as in \cite[Lemma~7.7]{A8}.  Here is the argument in the
present setting.  Both conventions use the same spin orientation
of the capped disk, the same complex orientation on sphere
factors, and the same reciprocal contractions.  They consequently
have no relative-class-dependent ratio.  At fixed index the
transverse graded-plane data are connected; frame paths with the
fixed spin data and determinant winding give the same conclusion.
The remaining factors arise from order, marking, and suspension
permutations and depend only on the corner indices.

To make the two-vertex argument explicit, use suspended degrees
$a_i=|x_i|-1$ and write the ratio as
$(-1)^{r_k(a_1,\ldots,a_k)}$, with $r_k$ valued in $\Z/2$.
The output suspended degree of a rigid vertex is
$1+\sum_i a_i$.  The determinant boundary identity says that
\begin{equation}
 \begin{split}
 &r_s(a_i,\ldots,a_{i+s-1})\\
 &\quad+r_{k-s+1}\left(a_1,\ldots,a_{i-1},
     1+\sum_{j=i}^{i+s-1}a_j,a_{i+s},\ldots,a_k\right)
 \end{split}
 \label{comp:sign-cocycle}
\end{equation}
is independent of the inserted block and its length.  Both
conventions have the same suspended insertion sign.  Thus
Equation~\eqref{comp:sign-cocycle} compares the same oriented
one-dimensional formal gluing problem.  It applies to arbitrary
integer cap indices, without assuming that the associated
holomorphic moduli space is nonempty.

Choose $h:\Z\to\Z/2$ recursively so that
$h(a+1)+h(a)=r_1(a)$.  Rescaling the line of suspended degree $a$
by $(-1)^{h(a)}$ replaces $r_k$ by
\[
 r'_k=r_k+h\left(1+\sum_i a_i\right)+\sum_i h(a_i).
\]
It makes $r'_1=0$ and preserves
Equation~\eqref{comp:sign-cocycle}.  Comparing a unary child at
input $i$ with a unary parent now gives
$r'_k(\ldots,a_i+1,\ldots)=r'_k(\ldots,a_i,\ldots)$.
Thus $r'_k$ is a constant $b_k$.  Comparing a binary child with a
unary parent gives $b_k=b_2+b_{k-1}$, so $b_k=(k-1)b_2$.
A common line rescaling by $(-1)^{b_2}$ removes these constants.
This proves a strict sign identification for both the ordinary
counts and the directed disk-integration comparison.  The
integer gradings and even factor dimensions used here are exactly
the features needed in this argument.

\subsection{The operadic diagonal and the product comparison}
\label{comp:operadic-completion}

Apply the count map $\gamma$ on $\mathcal G_0$ to the homotopy
chosen in Section~\ref{comp:domain-chains}.  At the first
endpoint, counts on the actual diagonal are the split product
polygon counts.  At the second endpoint, factorization on cross
products gives precisely the tensor-product operations determined
by $\Delta$.  For strips the two endpoints have the same tensor
differential: a rigid quotient product strip consists of a rigid
quotient strip in one factor and a constant strip in the other;
if both factors move, the independent translation gives excess
dimension.

The construction applies with object labels.  The endomorphism
operad simply has a component for each composable ordered word,
and its compositions use the corresponding slots.  Only the
bounded arities allowed by the finite strict ordering occur.  The
operadic path therefore gives isomorphic directed $A_\infty$
structures: integrate its coderivation homotopy successively in
arity over the characteristic-zero field.  The differential is
fixed and there is no arity-one homotopy term, so every step is a
finite integration of polynomial coefficients, with identity
linear term.

To finish, the cellular-diagonal tensor computes the ordinary
tensor through dg models and respects $A_\infty$
quasi-isomorphisms.  One may first forget units.  Bar--cobar
rectification replaces an $A_\infty$ quasi-isomorphism by a
zigzag of strict quasi-isomorphisms of operad algebras
\cite{LodayVallette}.  For fixed-object categories the same
construction uses composable tensor words over the finite
product of ground fields given by the object idempotents.
Units may be adjoined freely when augmented resolutions are
used.  Tensor these strict diagrams with
Equation~\eqref{comp:cellular-diagonal}; over a field their maps
remain quasi-isomorphisms.  At the dg endpoints only the
associative multiplication survives, and the resulting tensor is
the ordinary dg tensor.  This also explains the agreement with
the tensor-product framework in the product theorem of
\cite{AmorimProduct}.

Restricting these comparisons to the product labels and strict
words preserves every operation and functor equation used here.
The ordinary factor mirror equivalences then give the stated
directed comparison with their dg mirror categories.  This proves
the second assertion and completes
Proposition~\ref{prop:directed-comparison}.
